\documentclass{ximera}

\begin{document}
	\author{Stitz-Zeager}
	\xmtitle{Exercises for Graphs of Other Circular Functions}{}

\mfpicnumber{1} \opengraphsfile{ExercisesforGraphsofOtherCircularFunctions} % mfpic settings added 

\begin{question}
    
In Exercises \ref{othergraphsfirst} - \ref{othergraphslast}, graph one cycle of the given function.  State the period of the function.

\begin{problem}\label{othergraphsfirst}
    $y = \tan \left(t - \frac{\pi}{3} \right)$
\end{problem}

\begin{problem}
    $y = 2\tan \left( \frac{1}{4}t \right) - 3$
\end{problem}

\begin{problem}
     $y = \frac{1}{3}\tan(-2t - \pi) + 1$
\end{problem}

\begin{problem}
    $y = \sec \left( t - \frac{\pi}{2} \right)$
\end{problem}

\begin{problem}
    $y = -\csc \left( t + \frac{\pi}{3} \right)$ 
\end{problem}

\begin{problem}
    $y = -\frac{1}{3} \sec \left( \frac{1}{2}t + \frac{\pi}{3} \right)$
\end{problem}

\begin{problem}
    $y = \csc (2t - \pi)$
\end{problem}

\begin{problem}
    $y = \sec(3t - 2\pi) + 4$ 
\end{problem}

\begin{problem}
    $y = \csc \left( -t - \frac{\pi}{4} \right) - 2$
\end{problem}

\begin{problem}
    $y = \cot \left( t + \frac{\pi}{6} \right)$
\end{problem}

\begin{problem}
    $y = -11\cot \left( \frac{1}{5} t \right)$
\end{problem}

\begin{problem}\label{othergraphslast}
    $y = \frac{1}{3} \cot \left( 2t + \frac{3\pi}{2} \right) + 1$
\end{problem}

\end{question}

\begin{question}
    
In Exercises \ref{fitsecantcosecantfirst} - \ref{fitsecantcosecantlast},  the graph of a (co)secant function is given. Find a formula for the function in the form $F(t) = A \sec(\omega t + \phi) + B$ and $G(t) = A \csc(\omega t + \phi) + B$.  Select $\omega$ so  $\omega > 0$. Check your answer by graphing.

\begin{problem}
    

 Asymptotes:  $t = \pm \frac{\pi}{2}$, $t=\pm \frac{3\pi}{2}$, \dots \label{fitsecantcosecantfirst}  %$F(t) = 2 \sec(t-\pi)$, $G(t) = 2 \csc \left(t - \frac{\pi}{2} \right)$

\begin{mfpic}[13]{-6}{6}{-5}{5}
\tlabel[cc](6,-0.30){\scriptsize $t$}
\tlabel[cc](0.25,5){\scriptsize $y$}
\axes
\xmarks{-5 step 1 until 5}
\ymarks{-4, -3, 1, 2, 3, 4}
\point[4pt]{(-3.14159, 2), (0,-2), (3.14159,2)}
\tlabel[cc](-3.14159, 1.25){ \scriptsize $(-\pi, 2)$}
\gclear \tlabelrect(0, -1.25){ \scriptsize $(0,-2)$}
\tlabel[cc](3.14159,1.25){ \scriptsize $(\pi,2)$}
\tlpointsep{4pt}
%\axislabels {x}{{$\frac{\pi}{2}$} 1.5708, {$\pi$} 3.1416, {$\frac{3\pi}{2}$} 4.7124, {$2\pi$} 6.2832}
%\axislabels {y}{{$-3$} -3, {$3$} 3}
\dashed \polyline{(-4.712, -5), (-4.712, 5)}
\dashed \polyline{(4.712, -5), (4.712, 5)}
\dashed \polyline{(-1.57, -5), (-1.57, 5)}
\dashed \polyline{(1.57, -5), (1.57, 5)}
\penwd{1.25pt}
\arrow \reverse \arrow \function{-4.28, -2, 0.1}{2/(cos(x-pi))}
\arrow \reverse \arrow \function{-1.15, 1.15, 0.1}{2/(cos(x-pi))}
\arrow \reverse \arrow \function{2, 4.28, 0.1}{2/(cos(x-pi))}
\end{mfpic} 

\vfill

\columnbreak

\end{problem}

\begin{problem}
    
  Asymptotes:  $t = \pm 1$, $t = \pm 3$, $t = \pm 5$, \ldots  \label{fitsecantcosecantlast}  %$F(t) = \sec\left( \frac{\pi}{2} t \right) + 1$, $G(t) = \csc\left( \frac{\pi}{2} t + \frac{\pi}{2} \right) + 1$

\begin{mfpic}[18][13]{-5.5}{5.5}{-5}{5}
\tlabel[cc](5.5,-0.30){\scriptsize $t$}
\tlabel[cc](0.25,5){\scriptsize $y$}
\axes
\xmarks{-5 step 1 until 5}
\ymarks{-4, -3, -2, -1, 3, 4}
\point[4pt]{(-4,2), (-2,0), (0,2), (2,0), (4,2)}
\tlabel[cc](-4, 1.25){ \scriptsize $(-4, 2)$}
\tlabel[cc](-2, 0.75){ \scriptsize $(-2,0)$}
\gclear \tlabelrect(0, 1.25){ \scriptsize $(0,2)$}
\tlabel[cc](2, 0.75){ \scriptsize $(2,0)$}
\tlabel[cc](4, 1.25){ \scriptsize $(4, 2)$}
\tlpointsep{4pt}
%\axislabels {x}{{$\frac{\pi}{2}$} 1.5708, {$\pi$} 3.1416, {$\frac{3\pi}{2}$} 4.7124, {$2\pi$} 6.2832}
%\axislabels {y}{{$-3$} -3, {$3$} 3}
\dashed \polyline{(-5, -5), (-5, 5)}
\dashed \polyline{(5, -5), (5, 5)}
\dashed \polyline{(-3, -5), (-3, 5)}
\dashed \polyline{(3, -5), (3, 5)}
\dashed \polyline{(-1, -5), (-1, 5)}
\dashed \polyline{(1, -5), (1, 5)}
\penwd{1.25pt}
\arrow \reverse \arrow \function{-4.8, -3.2, 0.1}{1+(1/(cos(x*pi/2)))}
\arrow \reverse \arrow \function{-2.8, -1.2, 0.1}{1+(1/(cos(x*pi/2)))}
\arrow \reverse \arrow \function{-0.8, 0.8, 0.1}{1+(1/(cos(x*pi/2)))}
\arrow \reverse \arrow \function{1.2, 2.8,  0.1}{1+(1/(cos(x*pi/2)))}
\arrow \reverse \arrow \function{3.2, 4.8,  0.1}{1+(1/(cos(x*pi/2)))}
\end{mfpic} 

\end{problem}

\end{question}

\begin{question}
    
In Exercises \ref{fittangentfirst} - \ref{fittangentlast},  the graph of a (co)tangent function given. Find a formula the function in the form  $J(t) = A \tan(\omega t + \phi) + B$ and $K(t) = A \cot(\omega t + \phi) + B$.  Select $\omega$ so  $\omega > 0$.  Check your answer by graphing.

\begin{problem}

Asymptotes:  $t=-\frac{3 \pi}{4}$, $t=\frac{\pi}{4}$, $t = \frac{5\pi}{4}$, \dots \label{fittangentfirst}  %$J(t) = -\tan\left(t+ \frac{\pi}{4} \right)$, $K(t) = \cot  \left(t - \frac{\pi}{4} \right)$

\begin{mfpic}[18][13]{-4}{5.5}{-5}{5}
\tlabel[cc](5.5,-0.30){\scriptsize $t$}
\tlabel[cc](0.25,5){\scriptsize $y$}
\axes
\xmarks{-3 step 1 until 4}
\ymarks{-4 step 1 until 4}
\point[4pt]{(-0.7854, 0), (2.356,0), (0,-1), (1.57,1)}
\gclear \tlabelrect(-0.25, 0.75){ \scriptsize $\left(-\frac{\pi}{4}, 0 \right)$ \vphantom{$\frac{\pi}{4}$}}
\tlabel[cc](2, -1){ \scriptsize $\left(\frac{3\pi}{4}, 0 \right)$}
\tlabel[cc](2.5, 1.25){ \scriptsize $\left(\frac{\pi}{2}, 1 \right)$}
\tlabel[cc](-1.1,-1.25){\scriptsize $(0,-1)$}
\tlpointsep{4pt}
%\axislabels {x}{{$\frac{\pi}{2}$} 1.5708, {$\pi$} 3.1416, {$\frac{3\pi}{2}$} 4.7124, {$2\pi$} 6.2832}
%\axislabels {y}{{$-3$} -3, {$3$} 3}
\dashed \polyline{(-2.356, -5), (-2.356, 5)}
\dashed \polyline{(0.7854, -5), (0.7854, 5)}
\dashed \polyline{(3.92, -5), (3.92, 5)}
\penwd{1.25pt}
\arrow \reverse \arrow \function{-2.15, 0.558, 0.1}{-tan(x + (pi/4))}
\arrow \reverse \arrow \function{1, 3.7, 0.1}{-tan(x + (pi/4))}
\end{mfpic} 

\vfill

\columnbreak

\end{problem}

\begin{problem}
    
  Asymptotes:  $t = \pm 2$, $t = \pm 6$, $t = \pm 10$, \ldots  \label{fittangentlast}  %$J(t) = \tan\left( \frac{\pi}{4} t \right) + 1$, $K(t) = -\cot\left( \frac{\pi}{4} t + \frac{\pi}{2} \right) + 1$

\begin{mfpic}[13]{-6.5}{6.5}{-5}{5}
\tlabel[cc](6.5,-0.30){\scriptsize $t$}
\tlabel[cc](0.25,5){\scriptsize $y$}
\axes
\xmarks{-5 step 1 until 5}
\ymarks{-4, -3, -2, -1, 3, 4,2,1}
\point[4pt]{(-5,0), (0,1), (5,2)}
\tlabel[cc](-4, -0.75){ \scriptsize $(-5, 0)$}
\tlabel[cc](4, 2){ \scriptsize $(5,2)$}
\gclear \tlabelrect(0, 2){ \scriptsize $(0,1)$ \vphantom{$\
frac{3}{5}$}}
%\tlabel[cc](2, 0.5){ \scriptsize $(2,0)$}
%\tlabel[cc](4, 1.5){ \scriptsize $(4, 2)$}
\tlpointsep{4pt}
%\axislabels {x}{{$\frac{\pi}{2}$} 1.5708, {$\pi$} 3.1416, {$\frac{3\pi}{2}$} 4.7124, {$2\pi$} 6.2832}
%\axislabels {y}{{$-3$} -3, {$3$} 3}
\dashed \polyline{(-2, -5), (-2, 5)}
\dashed \polyline{(2, -5), (2, 5)}
\penwd{1.25pt}
\arrow \reverse \arrow \function{-5.7, -2.29, 0.1}{1+tan((pi*x)/4)}
\arrow \reverse \arrow \function{-1.7, 1.69, 0.1}{1+tan((pi*x)/4)}
\arrow \reverse \arrow \function{2.3, 5.69, 0.1}{1+tan((pi*x)/4)}
\end{mfpic}

\end{problem}

\end{question}


\begin{question}
In Section \ref{SineCosineLimits}, we observed\footnote{insisted?} that the cosine and sine functions are continuous.  As such, the other four circular functions are continuous on their domains.\footnote{See the remarks following Definition \ref{continuousdefn} in Section \ref{LimitPropertiesandContinuity}.} 
\smallskip
In Exercises \ref{othertriglimitexfirst} - \ref{othertriglimitexlast}, determine the given limit.  Use the symbols `$-\infty$' and `$\infty$'as appropriate.  Check your answers graphically.

\begin{problem}\label{othertriglimitexfirst} 
    $\ds{\lim_{t \rightarrow 0} \tan(t)}$.
\end{problem}

\begin{problem}
    $\ds{\lim_{t \rightarrow \pi} \sec(t)}$
\end{problem}

\begin{problem}
    $\ds{\lim_{t \rightarrow 3\pi}}$ $\cot\left(\frac{t}{2}\right)$
\end{problem}

\begin{problem}
     $\ds{\lim_{\theta \rightarrow 0}}$ $\csc\left(2\theta + \frac{\pi}{4}\right)$.
\end{problem}

\begin{problem}
    $\ds{\lim_{\theta \rightarrow \pi} (\cos(\theta)  -  \sec(\theta))}$
\end{problem}

\begin{problem}
    $\ds{\lim_{\theta \rightarrow \frac{\pi}{4}} \sec(\theta) \, \tan(\theta)}$
\end{problem}

\begin{problem}
    $\ds{\lim_{x \rightarrow 0^{+}} \csc(3x)}$
\end{problem}

\begin{problem}
    $\ds{\lim_{x \rightarrow \pi^{-}}  (\sin(2x)  - \tan(x))}$
\end{problem}

\begin{problem}\label{othertriglimitexlast} 
    $\ds{\lim_{x \rightarrow \pi^{+}}}$ $(\cos(x) + \cot(x))$
\end{problem}

\end{question}


\begin{problem}

In Exercise \ref{sintovertexercise3} in Section \ref{TheOtherCircularFunctions}, we proved $\ds{\lim_{\theta \rightarrow 0}}$ $\frac{\sin(\theta)}{\theta} = 1$.

\begin{enumerate}
\setcounter{enumi}{\value{HW}}

\item \label{thetaoversinetheta} Use Theorem \ref{LimitProp01} from Section \ref{IntroLimits} to show  $\ds{\lim_{\theta \rightarrow 0}}$ $\frac{\theta}{\sin(\theta)} = 1$.

\smallskip


\item  Which indeterminate form is present in the limit $\ds{\lim_{\theta \rightarrow 0^{+}} 2 \theta \, \csc(\theta)}$?

\smallskip

\item  Graph $f(\theta) = 2 \theta \, \csc(\theta)$ near $\theta = 0$.  What appears to be the limit?

\smallskip

\item  Rewrite $2 \theta \, \csc(\theta)$ in terms of $\theta$ and $\sin(\theta)$ and use part \ref{thetaoversinetheta} to analytically find $\ds{\lim_{\theta \rightarrow 0^{+}} 2 \theta \, \csc(\theta)}$.

\item  Use the same methodology as above  to help you analytically determine $\ds{\lim_{\theta \rightarrow 0^{+}} 2 \theta \, \cot(\theta)}$


\setcounter{HW}{\value{enumi}}
\end{enumerate}

\end{problem}


\begin{problem}\label{secantcosecantshiftexercise}  
    Use the conversion formulas listed in Theorem \ref{cosinesinefunctionprops} to create conversion formulas between secant and cosecant functions.
\end{problem}

\begin{problem}
    Use a conversion formula to rewrite our first answer to Example \ref{secantcosecantfromgraphex}, $f(t) = \sec\left( 2 t - \frac{7\pi}{6} \right) -1$, in terms of cosecants.
\end{problem}


\begin{problem}\label{NegativeAsecantcosecant} 
    Rework Example \ref{secantcosecantfromgraphex} and find answers with $A<0$.
\end{problem}

\begin{problem}\label{sinusoidformseccscexercise}  
    Prove Theorem \ref{secantcosecanttperiodphaseshift}    using Theorem  \ref{sinusoidform}.
\end{problem}

\begin{problem}\label{sinusoidformtancotexercise}
    Prove Theorem \ref{tangentcotangentperiodphaseshift} using Theorem \ref{transformationsthm}. 
\end{problem}

\begin{problem}\label{rangeoftangentexercise}
    
  In this Exercise, we argue the range of the tangent function is $(-\infty, \infty)$.  Let $M$ be a fixed, but arbitrary positive real number.

\begin{enumerate}

\item Show there is an acute  angle $\theta$ with $\tan(\theta)  = M$.  (Hint: think right triangles.) 

\item  Using the symmetry of the Unit Circle, explain why there are angles $\theta$ with $\tan(\theta) = -M$.

\item  Find angles with $\tan(\theta) = 0$.

\item  Combine the three parts above to conclude the range of the tangent function is $(-\infty, \infty)$.

\end{enumerate}

\end{problem}

\begin{problem}\label{cotisoddexercise}
    Prove $\cot(t)$ is odd.  (Hint:  mimic the proof given in the text that $\tan(t)$ is odd.)
\end{problem}
  

\end{document}
